
% Chapter 6

\chapter{Conclusion and Future Work} % Main chapter title

\label{Chapter6} % For referencing the chapter elsewhere, use \ref{Chapter6}

%-------------------------------------------------------------------------------

\section{Summary of findings}

This work presented a reproducible pipeline for constructing a multi-factor index that integrates ESG, financial and market signals, and for evaluating that index under a set of time-split protocols. The main practical outputs are:

\begin{itemize}
	\item A scoring and ranking pipeline that produces composite rankings with configurable weights.
	\item A weight-sensitivity grid search and visualisations that identify regions of robust ranking behaviour.
	\item A time-split evaluation framework (fixed, rolling, expanding) validated on synthetic panel data and used to produce the CSVs and figures referenced in the Results chapter.
\end{itemize}

These artifacts are included in the `Thesis_report/` folder so the Results narrative is directly traceable to the CSV and image outputs.

%-------------------------------------------------------------------------------

\section{Limitations}

The current study is intentionally conservative in scope so that the pipeline remains transparent and reproducible. Important limitations include:

\begin{itemize}
	\item The evaluation in this repository includes synthetic-panel validation; a full historical panel backtest requires time-indexed company-level histories that are not included here.
	\item Transaction costs and turnover impact have not been modelled in depth; performance metrics reported are gross returns unless specified in the CSVs.
	\item The scoring uses a fixed set of factor transformations and a pre-specified grid of weights; more flexible machine-learning based fusion methods were intentionally excluded to preserve interpretability.
\end{itemize}

%-------------------------------------------------------------------------------

\section{Future work}

We recommend the following practical extensions (ranked by suggested priority):

\begin{enumerate}
	\item \textbf{Run the pipeline on a full historical panel:} obtain time-stamped company-level data covering multiple years and re-run the time-split evaluation to produce fully historical out-of-sample metrics.
	\item \textbf{Incorporate transaction costs and turnover constraints:} simulate realistic rebalancing with explicit cost assumptions and report net returns and turnover statistics.
	\item \textbf{Explore rebalancing frequency and portfolio size:} evaluate sensitivity to rebalancing (monthly vs quarterly) and to the choice of top-N (e.g., top-10, top-20, top-50) to assess robustness.
	\item \textbf{Extend risk modelling:} integrate alternative risk estimates (e.g., factor models, shrinkage covariance estimators) to construct mean-variance efficient variants and to better contextualise the efficient frontier results.
	\item \textbf{Broaden the universes and ESG sources:} test robustness across larger universes and alternative ESG providers to assess data-source sensitivity.
	\item \textbf{Automation and reproducible builds:} provide a containerised TeX build and a short Makefile or script that assembles the `Thesis_report/` PDF from the included tables and figures.
\end{enumerate}

%-------------------------------------------------------------------------------

\section{Concluding remarks}

The project delivers a practical, auditable pipeline and a self-contained thesis package that demonstrates how ESG and financial signals can be combined and evaluated in a reproducible way. The next natural step is to run the same pipeline on full historical panels and to include turnover and cost adjustments so the results are ready for production decision-making.


